\chapter*{Annexe B}
\addcontentsline{toc}{chapter}{Annexe B --- Documentation de l'API REST}

\noindent Cette annexe documente l'API REST de SmartTrans-AI (préfixe \texttt{/api/v1}). La description interactive est également disponible via OpenAPI/Swagger : \url{http://127.0.0.1:8080/swagger-ui.html}.

\section*{B.1 Conventions}

\begin{itemize}
  \item Format d'échange : JSON (\texttt{Content-Type: application/json}).
  \item Identifiants : UUID.
  \item Authentification : en-tête \texttt{Authorization: Bearer <jeton>}, sauf \texttt{POST /api/v1/auth/login}.
  \item Rôles : \texttt{ADMIN}, \texttt{TRANSPORT\_MANAGER}, \texttt{DRIVER}.
\end{itemize}

\section*{B.2 Authentification}

\begin{table}[htbp]
\centering
\caption{Authentification}
\label{tab:api-auth}
\begin{tabular}{|p{0.38\textwidth}|p{0.14\textwidth}|p{0.38\textwidth}|}
\hline
\TableHead
\tht{Ressource} & \tht{Méthode} & \tht{Description} \\
\hline
\texttt{/api/v1/auth/login} & POST & Connexion ; retourne le JWT et le profil \\
\hline
\end{tabular}
\end{table}

\section*{B.3 Utilisateurs et rôles}

Les ressources suivantes sont réservées à l'administrateur.

\begin{table}[htbp]
\centering
\caption{Utilisateurs et rôles}
\label{tab:api-users}
\begin{tabular}{|p{0.38\textwidth}|p{0.14\textwidth}|p{0.38\textwidth}|}
\hline
\TableHead
\tht{Ressource} & \tht{Méthode} & \tht{Description} \\
\hline
\texttt{/api/v1/users} & GET, POST & Liste et création des comptes \\
\hline
\texttt{/api/v1/users/\{id\}} & GET, PATCH, DELETE & Lecture, mise à jour, suppression \\
\hline
\texttt{/api/v1/roles} & GET, POST & Liste et création des rôles \\
\hline
\texttt{/api/v1/roles/\{id\}} & GET, PATCH, DELETE & Lecture, mise à jour, suppression \\
\hline
\end{tabular}
\end{table}

\section*{B.4 Référentiels métier}

Clients, chauffeurs et véhicules sont gérés par l'administrateur et le responsable transport. Le chauffeur accède à son profil (\texttt{/drivers/me}), à la consultation limitée des fiches et à la mise à jour de sa disponibilité.

\begin{table}[htbp]
\centering
\caption{Clients, chauffeurs et véhicules}
\label{tab:api-ref}
\begin{tabular}{|p{0.44\textwidth}|p{0.14\textwidth}|p{0.32\textwidth}|}
\hline
\TableHead
\tht{Ressource} & \tht{Méthode} & \tht{Description} \\
\hline
\texttt{/api/v1/clients} & GET, POST & Liste et création \\
\hline
\texttt{/api/v1/clients/\{id\}} & GET, PATCH, DELETE & Fiche client \\
\hline
\texttt{/api/v1/drivers} & GET, POST & Liste et création \\
\hline
\texttt{/api/v1/drivers/me} & GET & Profil du chauffeur connecté \\
\hline
\texttt{/api/v1/drivers/expiring-licenses} & GET & Permis proches de l'échéance \\
\hline
\texttt{/api/v1/drivers/\{id\}} & GET, PATCH, DELETE & Fiche chauffeur \\
\hline
\texttt{/api/v1/drivers/\{id\}/availability} & PATCH & Disponibilité \\
\hline
\texttt{/api/v1/drivers/\{id\}/deactivate} & PATCH & Désactivation \\
\hline
\texttt{/api/v1/vehicles} & GET, POST & Liste et création \\
\hline
\texttt{/api/v1/vehicles/expiring-documents} & GET & Assurance / visite technique \\
\hline
\texttt{/api/v1/vehicles/\{id\}} & GET, PATCH, DELETE & Fiche véhicule \\
\hline
\texttt{/api/v1/vehicles/\{id\}/deactivate} & PATCH & Désactivation \\
\hline
\end{tabular}
\end{table}

\section*{B.5 Missions}

Le chauffeur peut consulter les missions qui le concernent et enchaîner les statuts autorisés (planifiée $\rightarrow$ en cours $\rightarrow$ terminée). L'annulation reste du ressort de l'exploitation.

\begin{table}[htbp]
\centering
\caption{Missions}
\label{tab:api-missions}
\begin{tabular}{|p{0.42\textwidth}|p{0.14\textwidth}|p{0.34\textwidth}|}
\hline
\TableHead
\tht{Ressource} & \tht{Méthode} & \tht{Description} \\
\hline
\texttt{/api/v1/missions} & GET, POST & Liste et planification \\
\hline
\texttt{/api/v1/missions/\{id\}} & GET, PATCH, DELETE & Fiche mission \\
\hline
\texttt{/api/v1/missions/\{id\}/status} & PATCH & Changement de statut \\
\hline
\end{tabular}
\end{table}

\section*{B.6 Journal et assistant}

Le journal est en lecture seule. L'assistant est réservé à l'administrateur et au responsable transport.

\begin{table}[htbp]
\centering
\caption{Journal et assistant}
\label{tab:api-pilotage}
\begin{tabular}{|p{0.40\textwidth}|p{0.14\textwidth}|p{0.36\textwidth}|}
\hline
\TableHead
\tht{Ressource} & \tht{Méthode} & \tht{Description} \\
\hline
\texttt{/api/v1/activities} & GET & Journal d'audit \\
\hline
\texttt{/api/v1/activities/\{id\}} & GET & Détail d'un événement \\
\hline
\texttt{/api/v1/assistant/ask} & POST & Question métier ancrée sur les données \\
\hline
\end{tabular}
\end{table}

\section*{B.7 Accès à la documentation interactive}

Avec le backend démarré (\texttt{./mvnw spring-boot:run}), ouvrir :

\url{http://127.0.0.1:8080/swagger-ui.html}

L'état de l'API est indiqué par \url{http://127.0.0.1:8080/actuator/health}.
